\documentclass[a4paper,10pt]{article}
\usepackage{fontspec}
\usepackage{xeCJK}
\usepackage[a4paper,margin=20mm]{geometry}
\usepackage{graphicx}
\usepackage{fancyhdr}
\setmonofont{Courier New}
\setCJKmainfont{Hiragino Sans W3}
\setCJKsansfont{Hiragino Sans W3}
\setCJKmonofont{Hiragino Sans W3}
\XeTeXlinebreaklocale ""
\xeCJKsetup{CJKglue={},CJKecglue={}}
\newdimen\charwd
{\tt\global\setbox0=\hbox{x}\global\charwd=\wd0}
\frenchspacing
\pagestyle{fancy}
\renewcommand{\headrulewidth}{0pt}
\fancyhf{}
\fancyhead[R]{\rm src2tex version 2.236}
\fancyfoot[R]{\rm input{\tt /}hanoi.swift\qquad page \thepage}
% plain TeX compatibility
\makeatletter
\providecommand{\eqalign}[1]{\vcenter{\openup1\jot\m@th
  \ialign{\strut\hfil$\displaystyle{##}$&$\displaystyle{{}##}$\hfil\crcr#1\crcr}}}
\makeatother
\ifx\sevenrm\undefined
  \font\sevenrm=cmr7 scaled \magstep0
\fi
\def\mc{\relax}
\def\gt{\relax}
\def\sc{\scshape}
\begin{document}
\tt\mc 

\noindent
\mbox{}\rm\mc {\tt /}{\tt *}\  \hrulefill *


\ % beginning of TeX mode

\centerline{\bf Towers of Hanoi (Swift)}

\begin{quote}

This program gives an answer to the following famous problem (towers of
Hanoi).
There is a legend that when one of the temples in Hanoi was constructed,
three poles were erected and a tower consisting of 64 golden discs was
arranged on one pole, their sizes decreasing regularly from bottom to top.
The monks were to move the tower of discs to the opposite pole, moving
only one at a time, and never putting any size disc above a smaller one.
The job was to be done in the minimum numbers of moves. What strategy for
moving discs will accomplish this optimum transfer?
\end{quote}

% end of TeX mode 

* \hrulefill \rm\mc \ {\tt *}{\tt /}
\tt\mc 

\noindent
\mbox{}\hfill

\noindent
\mbox{}\hfill

\noindent
\mbox{}\rm\mc {\tt /}{\tt *}\  \hrulefill\ hanoi.swift \ \hrulefill \rm\mc \ {\tt *}{\tt /}
\tt\mc 

\noindent
\mbox{}\hfill

\noindent
\mbox{}\hfill

\noindent
\mbox{}let{\tt\mc \ }ARRAY{\tt\mc \ }={\tt\mc \ }8{\tt\mc \ \ \ \ \ \ \ \ }{\tt\mc \ \ \ \ \ \ \ \ }{\tt\mc \ \ \ \ \ \ \ \ }{\tt\mc \ \ \ \ \ \ \ \ }\rm\mc {\tt /}{\tt *}\  \ disc の数 \hfill \rm\mc \ {\tt *}{\tt /}
\tt\mc 

\noindent
\mbox{}\hfill

\noindent
\mbox{}var{\tt\mc \ }disc{\tt\mc \ }={\tt\mc \ }[[Int]](repeating:{\tt\mc \ }[Int](repeating:{\tt\mc \ }0,{\tt\mc \ }count:{\tt\mc \ }8),{\tt\mc \ }count:{\tt\mc \ }3)

\noindent
\mbox{}{\tt\mc \ \ \ \ \ \ \ \ }{\tt\mc \ \ \ \ \ \ \ \ }{\tt\mc \ \ \ \ \ \ \ \ }{\tt\mc \ \ \ \ \ \ \ \ }{\tt\mc \ \ \ \ \ \ \ \ }\rm\mc {\tt /}{\tt *}\  \ disc に関するデータの置き場所
					   \hfill \rm\mc \ {\tt *}{\tt /}
\tt\mc 

\noindent
\mbox{}\hfill

\noindent
\mbox{}func{\tt\mc \ }initArray(){\tt\mc \ }{\tt\char'173}{\tt\mc \ \ \ \ \ \ \ \ }{\tt\mc \ \ \ \ \ \ \ \ }{\tt\mc \ \ \ \ \ \ \ \ }\rm\mc {\tt /}{\tt *}\  \ disc に関するデータの初期化
					   \hfill \rm\mc \ {\tt *}{\tt /}
\tt\mc 

\noindent
\mbox{}{\tt\mc \ }{\tt\mc \ }{\tt\mc \ }{\tt\mc \ }for{\tt\mc \ }j{\tt\mc \ }in{\tt\mc \ }0..{\tt <}ARRAY{\tt\mc \ }{\tt\char'173}

\noindent
\mbox{}{\tt\mc \ }{\tt\mc \ }{\tt\mc \ }{\tt\mc \ }{\tt\mc \ }{\tt\mc \ }{\tt\mc \ }{\tt\mc \ }disc[0][j]{\tt\mc \ }={\tt\mc \ }ARRAY{\tt\mc \ }{\tt -}{\tt\mc \ }j

\noindent
\mbox{}{\tt\mc \ }{\tt\mc \ }{\tt\mc \ }{\tt\mc \ }{\tt\mc \ }{\tt\mc \ }{\tt\mc \ }{\tt\mc \ }disc[1][j]{\tt\mc \ }={\tt\mc \ }0

\noindent
\mbox{}{\tt\mc \ }{\tt\mc \ }{\tt\mc \ }{\tt\mc \ }{\tt\mc \ }{\tt\mc \ }{\tt\mc \ }{\tt\mc \ }disc[2][j]{\tt\mc \ }={\tt\mc \ }0

\noindent
\mbox{}{\tt\mc \ }{\tt\mc \ }{\tt\mc \ }{\tt\mc \ }{\tt\char'175}

\noindent
\mbox{}{\tt\char'175}

\noindent
\mbox{}\hfill

\noindent
\mbox{}var{\tt\mc \ }counter{\tt\mc \ }={\tt\mc \ }0{\tt\mc \ \ \ \ \ \ \ \ }{\tt\mc \ \ \ \ \ \ \ \ }{\tt\mc \ \ \ \ \ \ \ \ }{\tt\mc \ \ \ \ \ \ \ \ }\rm\mc {\tt /}{\tt *}\  \ 移動回数カウンタ \hfill \rm\mc \ {\tt *}{\tt /}
\tt\mc 

\noindent
\mbox{}\hfill

\noindent
\mbox{}func{\tt\mc \ }printResult(){\tt\mc \ }{\tt\char'173}{\tt\mc \ \ \ \ \ \ \ \ }{\tt\mc \ \ \ \ \ \ \ \ }{\tt\mc \ \ \ \ \ \ \ \ }\rm\mc {\tt /}{\tt *}\  \ 結果の表示 \hfill \rm\mc \ {\tt *}{\tt /}
\tt\mc 

\noindent
\mbox{}{\tt\mc \ }{\tt\mc \ }{\tt\mc \ }{\tt\mc \ }counter{\tt\mc \ }+={\tt\mc \ }1

\noindent
\mbox{}{\tt\mc \ }{\tt\mc \ }{\tt\mc \ }{\tt\mc \ }print({\tt "}{\tt -}{\tt -}{\tt -}{\tt\#}{\tt\char92}(counter){\tt -}{\tt -}{\tt -}{\tt "})

\noindent
\mbox{}{\tt\mc \ }{\tt\mc \ }{\tt\mc \ }{\tt\mc \ }for{\tt\mc \ }i{\tt\mc \ }in{\tt\mc \ }0...2{\tt\mc \ }{\tt\char'173}

\noindent
\mbox{}{\tt\mc \ }{\tt\mc \ }{\tt\mc \ }{\tt\mc \ }{\tt\mc \ }{\tt\mc \ }{\tt\mc \ }{\tt\mc \ }print({\tt "}[{\tt\char92}(i)]{\tt\mc \ }{\tt "},{\tt\mc \ }terminator:{\tt\mc \ }{\tt "}{\tt "})

\noindent
\mbox{}{\tt\mc \ }{\tt\mc \ }{\tt\mc \ }{\tt\mc \ }{\tt\mc \ }{\tt\mc \ }{\tt\mc \ }{\tt\mc \ }for{\tt\mc \ }j{\tt\mc \ }in{\tt\mc \ }0..{\tt <}ARRAY{\tt\mc \ }{\tt\char'173}

\noindent
\mbox{}{\tt\mc \ }{\tt\mc \ }{\tt\mc \ }{\tt\mc \ }{\tt\mc \ }{\tt\mc \ }{\tt\mc \ }{\tt\mc \ }{\tt\mc \ }{\tt\mc \ }{\tt\mc \ }{\tt\mc \ }if{\tt\mc \ }disc[i][j]{\tt\mc \ }!={\tt\mc \ }0{\tt\mc \ }{\tt\char'173}

\noindent
\mbox{}{\tt\mc \ }{\tt\mc \ }{\tt\mc \ }{\tt\mc \ }{\tt\mc \ }{\tt\mc \ }{\tt\mc \ }{\tt\mc \ }{\tt\mc \ }{\tt\mc \ }{\tt\mc \ }{\tt\mc \ }{\tt\mc \ }{\tt\mc \ }{\tt\mc \ }{\tt\mc \ }print({\tt "}{\tt\char92}(disc[i][j]){\tt\mc \ }{\tt "},{\tt\mc \ }terminator:{\tt\mc \ }{\tt "}{\tt "})

\noindent
\mbox{}{\tt\mc \ }{\tt\mc \ }{\tt\mc \ }{\tt\mc \ }{\tt\mc \ }{\tt\mc \ }{\tt\mc \ }{\tt\mc \ }{\tt\mc \ }{\tt\mc \ }{\tt\mc \ }{\tt\mc \ }{\tt\char'175}{\tt\mc \ }else{\tt\mc \ }{\tt\char'173}

\noindent
\mbox{}{\tt\mc \ }{\tt\mc \ }{\tt\mc \ }{\tt\mc \ }{\tt\mc \ }{\tt\mc \ }{\tt\mc \ }{\tt\mc \ }{\tt\mc \ }{\tt\mc \ }{\tt\mc \ }{\tt\mc \ }{\tt\mc \ }{\tt\mc \ }{\tt\mc \ }{\tt\mc \ }break

\noindent
\mbox{}{\tt\mc \ }{\tt\mc \ }{\tt\mc \ }{\tt\mc \ }{\tt\mc \ }{\tt\mc \ }{\tt\mc \ }{\tt\mc \ }{\tt\mc \ }{\tt\mc \ }{\tt\mc \ }{\tt\mc \ }{\tt\char'175}

\noindent
\mbox{}{\tt\mc \ }{\tt\mc \ }{\tt\mc \ }{\tt\mc \ }{\tt\mc \ }{\tt\mc \ }{\tt\mc \ }{\tt\mc \ }{\tt\char'175}

\noindent
\mbox{}{\tt\mc \ }{\tt\mc \ }{\tt\mc \ }{\tt\mc \ }{\tt\mc \ }{\tt\mc \ }{\tt\mc \ }{\tt\mc \ }print()

\noindent
\mbox{}{\tt\mc \ }{\tt\mc \ }{\tt\mc \ }{\tt\mc \ }{\tt\char'175}

\noindent
\mbox{}{\tt\char'175}

\noindent
\mbox{}\hfill

\noindent
\mbox{}var{\tt\mc \ }ptr{\tt\mc \ }={\tt\mc \ }[0,{\tt\mc \ }0,{\tt\mc \ }0]{\tt\mc \ \ \ \ \ \ \ \ }{\tt\mc \ \ \ \ \ \ \ \ }{\tt\mc \ \ \ \ \ \ \ \ }\rm\mc {\tt /}{\tt *}\  \ disc 移動用ポインタ（インデックス）
					   \hfill \rm\mc \ {\tt *}{\tt /}
\tt\mc 

\noindent
\mbox{}\hfill

\noindent
\mbox{}func{\tt\mc \ }moveOneDisc({\tt\_}{\tt\mc \ }i:{\tt\mc \ }Int,{\tt\mc \ }{\tt\_}{\tt\mc \ }j:{\tt\mc \ }Int){\tt\mc \ }{\tt\char'173}{\tt\mc \ \ \ \ \ \ \ \ }\rm\mc {\tt /}{\tt *}\  \ 1枚の disc を pole $i$ から
					   pole $j$ に移動する \hfill \rm\mc \ {\tt *}{\tt /}
\tt\mc 

\noindent
\mbox{}{\tt\mc \ }{\tt\mc \ }{\tt\mc \ }{\tt\mc \ }ptr[i]{\tt\mc \ }{\tt -}={\tt\mc \ }1

\noindent
\mbox{}{\tt\mc \ }{\tt\mc \ }{\tt\mc \ }{\tt\mc \ }disc[j][ptr[j]]{\tt\mc \ }={\tt\mc \ }disc[i][ptr[i]]

\noindent
\mbox{}{\tt\mc \ }{\tt\mc \ }{\tt\mc \ }{\tt\mc \ }ptr[j]{\tt\mc \ }+={\tt\mc \ }1

\noindent
\mbox{}{\tt\mc \ }{\tt\mc \ }{\tt\mc \ }{\tt\mc \ }disc[i][ptr[i]]{\tt\mc \ }={\tt\mc \ }0

\noindent
\mbox{}{\tt\char'175}

\noindent
\mbox{}\hfill

\noindent
\mbox{}func{\tt\mc \ }moveDiscs({\tt\_}{\tt\mc \ }n:{\tt\mc \ }Int,{\tt\mc \ }{\tt\_}{\tt\mc \ }i:{\tt\mc \ }Int,{\tt\mc \ }{\tt\_}{\tt\mc \ }j:{\tt\mc \ }Int,{\tt\mc \ }{\tt\_}{\tt\mc \ }k:{\tt\mc \ }Int){\tt\mc \ }{\tt\char'173}

\noindent
\mbox{}{\tt\mc \ \ \ \ \ \ \ \ }{\tt\mc \ \ \ \ \ \ \ \ }{\tt\mc \ \ \ \ \ \ \ \ }{\tt\mc \ \ \ \ \ \ \ \ }{\tt\mc \ \ \ \ \ \ \ \ }\rm\mc {\tt /}{\tt *}\  \ 上から $n$ 枚目までの disc
					   を、pole $i$ から pole $j$ に
					   \hfill \rm\mc \ {\tt *}{\tt /}
\tt\mc 

\noindent
\mbox{}{\tt\mc \ \ \ \ \ \ \ \ }{\tt\mc \ \ \ \ \ \ \ \ }{\tt\mc \ \ \ \ \ \ \ \ }{\tt\mc \ \ \ \ \ \ \ \ }{\tt\mc \ \ \ \ \ \ \ \ }\rm\mc {\tt /}{\tt *}\  \ pole $k$ を経由して、移動する
					   \hfill \rm\mc \ {\tt *}{\tt /}
\tt\mc 

\noindent
\mbox{}{\tt\mc \ }{\tt\mc \ }{\tt\mc \ }{\tt\mc \ }if{\tt\mc \ }n{\tt\mc \ }{\tt >}={\tt\mc \ }1{\tt\mc \ }{\tt\char'173}

\noindent
\mbox{}{\tt\mc \ }{\tt\mc \ }{\tt\mc \ }{\tt\mc \ }{\tt\mc \ }{\tt\mc \ }{\tt\mc \ }{\tt\mc \ }moveDiscs(n{\tt\mc \ }{\tt -}{\tt\mc \ }1,{\tt\mc \ }i,{\tt\mc \ }k,{\tt\mc \ }j){\tt\mc \ \ \ \ \ \ \ \ }\rm\mc {\tt /}{\tt *}\  \ 関数 {\tt moveDiscs()}
					   の中で、さらに自分自身 \hfill \rm\mc \ {\tt *}{\tt /}
\tt\mc 

\noindent
\mbox{}{\tt\mc \ }{\tt\mc \ }{\tt\mc \ }{\tt\mc \ }{\tt\mc \ }{\tt\mc \ }{\tt\mc \ }{\tt\mc \ }moveOneDisc(i,{\tt\mc \ }j){\tt\mc \ \ \ \ \ \ \ \ }{\tt\mc \ \ \ \ \ \ \ \ }\rm\mc {\tt /}{\tt *}\  \ {\tt moveDiscs()} が使われ
					   ている。このような \hfill \rm\mc \ {\tt *}{\tt /}
\tt\mc 

\noindent
\mbox{}{\tt\mc \ }{\tt\mc \ }{\tt\mc \ }{\tt\mc \ }{\tt\mc \ }{\tt\mc \ }{\tt\mc \ }{\tt\mc \ }printResult(){\tt\mc \ \ \ \ \ \ \ \ }{\tt\mc \ \ \ \ \ \ \ \ }{\tt\mc \ \ \ \ \ \ \ \ }\rm\mc {\tt /}{\tt *}\  \ 手法は、「再帰的呼びだし」
					   といわれる。 \hfill \rm\mc \ {\tt *}{\tt /}
\tt\mc 

\noindent
\mbox{}{\tt\mc \ }{\tt\mc \ }{\tt\mc \ }{\tt\mc \ }{\tt\mc \ }{\tt\mc \ }{\tt\mc \ }{\tt\mc \ }moveDiscs(n{\tt\mc \ }{\tt -}{\tt\mc \ }1,{\tt\mc \ }k,{\tt\mc \ }j,{\tt\mc \ }i)

\noindent
\mbox{}{\tt\mc \ }{\tt\mc \ }{\tt\mc \ }{\tt\mc \ }{\tt\char'175}

\noindent
\mbox{}{\tt\char'175}

\noindent
\mbox{}\hfill

\noindent
\mbox{}\rm\mc {\tt /}{\tt *}\  \par\begin{center}

\includegraphics[scale=0.3]{hanoi1}\quad
\includegraphics[scale=0.3]{hanoi2}\end{center}


たとえば、関数 {\tt moveDiscs(4, 0, 1, 2)} を呼び出すことは、
上図のような操作をすることに対応する。\hfill \rm\mc \ {\tt *}{\tt /}
\tt\mc 

\noindent
\mbox{}\hfill

\noindent
\mbox{}\hfill

\noindent
\mbox{}ptr[0]{\tt\mc \ }={\tt\mc \ }ARRAY

\noindent
\mbox{}ptr[1]{\tt\mc \ }={\tt\mc \ }0

\noindent
\mbox{}ptr[2]{\tt\mc \ }={\tt\mc \ }0

\noindent
\mbox{}\hfill

\noindent
\mbox{}initArray()

\noindent
\mbox{}moveDiscs(ARRAY,{\tt\mc \ }0,{\tt\mc \ }1,{\tt\mc \ }2){\tt\mc \ \ \ \ \ \ \ \ }{\tt\mc \ \ \ \ \ \ \ \ }\rm\mc {\tt /}{\tt *}\  \ {\tt ARRAY} 枚の disc を
					   pole 0 から pole 1 に pole 2
					   \hfill \rm\mc \ {\tt *}{\tt /}
\tt\mc 

\noindent
\mbox{}{\tt\mc \ \ \ \ \ \ \ \ }{\tt\mc \ \ \ \ \ \ \ \ }{\tt\mc \ \ \ \ \ \ \ \ }{\tt\mc \ \ \ \ \ \ \ \ }{\tt\mc \ \ \ \ \ \ \ \ }\rm\mc {\tt /}{\tt *}\  \ を経由して、移動する \hfill \rm\mc \ {\tt *}{\tt /}

\tt\mc 

\rm\mc

\end{document}
